problem:  
Let \(G\) be a connected, simply connected **eight‑dimensional Lie group** endowed with a **left‑invariant Hermitian structure** \((J,g)\), where \(J\) is integrable and the Lee form \(\theta\) is exact (\(d\Omega^3=\theta\wedge\Omega^3\), \(\theta=d\varphi\)), so that \(\tilde g=e^{-\varphi}g\) is **balanced**.  

Suppose now that \(G\) admits a smooth, not necessarily left‑invariant, **\(J\)-invariant integrable distribution** \(D\subset TG\) of real rank \(4\), such that the induced foliation \(\mathcal{F}\) has **complex surfaces as leaves**, each **minimal** with respect to \(g\), and \(D\cap Z(\mathfrak g)=\{0\}\).  

**Conjectural rigidity statement:**  
Do the hypotheses of **integrability of \(J\), exactness of the Lee form, and minimality of \(\mathcal{F}\)** force \(D\) to be **locally left‑invariant** up to automorphism of the complex Lie structure—that is, locally generated by a complex Lie subalgebra of \(\mathfrak g\)?  

Equivalently, does every minimal holomorphic foliation on a globally conformally balanced eight‑dimensional Hermitian Lie group arise algebraically from the Lie structure itself?  

If not, can non‑invariant minimal foliations exist, and can possible obstructions to invariance be expressed via the (1,1)‑curvature component of the **Chern connection** or the non‑parallel part of the complex mean‑curvature form along \(D^\perp\)?  

Why is it a "good" problem:  
This refined problem eliminates the triviality of left‑invariance and pinpoints a truly open and subtle rigidity question: whether analytic and geometric conditions (integrable complex structure, conformal balance, and minimality) can *force* a foliation to become algebraic in nature.  
It lies at the interplay of **Hermitian torsion theory**, **foliation geometry**, and **Lie group algebraic structure**, bridging differential and algebraic viewpoints on complex geometry.  
A resolution would illuminate how analytic curvature data of the Chern connection constrain global symmetry, providing a deep link between conformal Hermitian analysis and Lie‑algebraic rigidity—a setting simple in statement yet complex in depth, fitting the profile of a “narrow entrance, profound depth” mathematical problem.